\documentclass{sopra-base}

\usepackage{sopra-documentation}
\usepackage{sopra-attachments}

\title{The 'sopra-attachments' Package}
\subtitle[Documentation for the 'sopra-attachments' package]{Documentation for the 'sopra-attachments' package | Version \thesoaversion}
\duedate{2020-01-5}

\keywords{Documentation,sopra-attachments,sopra,uni ulm,uulm,package}
\authors{Florian Sihler (florian.sihler@uni-ulm.de)}

\group{Die Affenbande}
\begin{document}
    \maketitle

\section{General}
\subsection{Why, what for?}
    This \LaTeXe\ package was written as part of the Sopra in the
    winter semester 2019 and summer semester 2020 and serves as a
    small wrapper for embedding documents
    of \imptext{Team 20}. This documentation was created together with the
    \T{sopra-base.cls} as well as the \T{sopra-documentation.sty} package.\par
    The \T{sopra-listings} package is used to visualize the individual code
    excerpts.
    The associated package should also be embedded in this document: \attachTexText{sopra-attachments.sty}{sopra-attachments.sty}.
\subsection{Dependencies}
    This package includes the following packages:
    \begin{multicols}{3}
        \begin{itemize}
            \def\pkgparse#1:#2\@nil{%
                \T{#1}\ifx!#2!\else\textsuperscript{(#2)}\fi%
            }
            \foreach \pkg in {xcolor:,attachfile2:,graphicx:{\jmark[\T{attsymbols}]{mrk:attsymbols}, \jmark[\T{usefa}]{mrk:usefa}},fontawesome:{\jmark[\T{usefa}]{mrk:usefa}}} {
                \item \expandafter\pkgparse\pkg\@nil
            }
        \end{itemize}
    \end{multicols}
    All of these packages should be part of a common \LaTeX\ distribution.

\subsection{Installation}
    The package is not delivered as a \T{.dtx}, which is why the
    following options arise:
    \begin{itemize}
        \item The package can be placed in the same directory as the
                document. In this case the inclusion instruction reads:
\begin{plainlatex}
\usepackage{sopra-attachments}
\end{plainlatex}
        \item The package can be placed in a subdirectory/a directory
                shipped with the document. In this case it is specified
                via the (relative) path:
\begin{plainlatex}
\usepackage{./My/Path/to/sopra-attachments}
\end{plainlatex}
        \item The package can be placed (via a symlink or similar)
              in one's own \emph{texmf} tree.
              For example, on Linux using texlive,
              the package can be placed here: \bvoid{\~/texmf/tex/latex/}.
              The directory can be created and then updated using
              \bbash{texhash \~/texmf}. Now
              the package can be used like any other installed package:
\begin{plainlatex}
\usepackage{sopra-attachments}
\end{plainlatex}
    \end{itemize}

\subsection{Further peculiarities}
In version \thesoaversion{} (\cmdref{thesoaversion}) the following applies: no default color is set anymore, to make it easier to configure it yourself. Furthermore \argref{usefa}{nousefa} was added.

\subsection{Accepted Parameters}
    Like most packages, this package accepts arguments.
    For arguments with a \say{counter} option, the one that is respectively active by default is written first and the other one in parentheses.
    So implicitly:
\begin{plainlatex}
    \usepackage[attsymbols,nodocheck,usefa]{sopra-attachments}
\end{plainlatex}
    is called. Whereas with:
\begin{plainlatex}
    \usepackage[docheck]{sopra-attachments}
\end{plainlatex}
    no error is thrown if the file does not exist!

    \begin{argument}{attsymbols}{noattsymbols}
        With \T{noattsymbols}, text is used instead even in cases where a symbol
        would actually be used.
    \end{argument}

    \begin{argument}{nodocheck}{docheck}
        With \T{docheck} it is checked whether the file actually exists. It is only attached
        if it exists; otherwise \emph{no} error is thrown. This can, for example,
        be useful when the graphic to be embedded is only generated later in the document
        using Ti\textit{k}Z's \T{external}.
    \end{argument}

    \begin{argument}{usefa}{nousefa}
        With \T{usefa}, the packages \emph{fontawesome} and \emph{graphicx} are loaded. A paperclip is prepended to the attached files for the text commands like \cmdref{attachDocumentText}, to make the attachment clear. With \T{nousefa} you get the behavior from before version \T{v1.1.00} back.
    \end{argument}

\subsection{Further Notes}

    Using \cmd{attachfilesetup} from the \T{attachfile2} package, further
    settings can be made. If \T{sopra-base} is not used, the author
    of the documents must be set via \cmd{author}.

\section{Commands and Environments}

\subsection{General Commands}

\begin{command}{thesoaversion}{}
    Returns the current version of the package. So \cmd{thesoaversion} results in: \thesoaversion\\
    \notetext{Note: via \blatex{\\value\{soaversion\}} the
    version can be obtained as a $4$-digit number: \arabic{soaversion}.}
\end{command}

\begin{command}{attachPdf}{\optArg{search path}\manArg{File}}
    Attempts to embed \T{File}. If it is not located in the same folder,
    \T{search path} can be specified as a prefix. \notetext{A trailing \say{slash} is
    expected!}
\end{command}

\begin{command}{attachPdfText}{\optArg{search path}\manArg{File}\manArg{Text}}
    Analogous to \cmdref{attachPdf}, but instead of the symbol (if enabled),
    the \T{Text} is displayed.
\end{command}

\begin{command}{attachDocument}{\optArg{search path}\manArg{File}}
    Analogous to \cmdref{attachPdf}, but the document type is set to \T{text/plain}. \notetext{The alias \cmd{attachTex} exists.}
\end{command}

\begin{command}{attachDocumentText}{\optArg{search path}\manArg{File}\manArg{Text}}
    Analogous to \cmdref{attachPdfText}, but the document type is set to \T{text/plain}. \notetext{The alias \cmd{attachTexText} exists.}
\end{command}

\end{document}
